\documentclass[11pt]{article}

\usepackage[margin=1in]{geometry}
\usepackage{amsmath, amssymb, amsthm}
\usepackage{mathtools}
\usepackage{bm}
\usepackage{xcolor}
\usepackage{hyperref}
\usepackage[most]{tcolorbox}
\usepackage{booktabs}
\usepackage{enumitem}
\usepackage{listings}

\hypersetup{colorlinks=true, linkcolor=blue!50!black, urlcolor=blue!50!black,
            unicode=true, bookmarksdepth=2}

\definecolor{codebg}{HTML}{F3F4F6}
\definecolor{accent}{HTML}{C85020}
\definecolor{navy}{HTML}{12264E}
\definecolor{ok}{HTML}{1F7A4C}

\lstdefinestyle{stata}{
  basicstyle=\ttfamily\small,
  backgroundcolor=\color{codebg},
  frame=none, breaklines=true, columns=fullflexible,
  keepspaces=true, xleftmargin=4pt, framexleftmargin=4pt,
  commentstyle=\color{gray!90!black}\itshape,
  keywordstyle=\color{navy}\bfseries,
  morekeywords={reghdfe, hdfe, gen, generate, replace, drop, capture,
                bys, bysort, gegen, levelsof, forvalues, foreach, if,
                local, scalar, matrix, predict, sum, total, accum, invsym,
                det, import, use, save, tempfile, syntax, varlist},
  morecomment=[l]{*},
  morecomment=[l]{//},
}
\lstset{style=stata}

\newcommand{\code}[1]{\texttt{\small #1}}
\newcommand{\Adoref}[1]{\textcolor{accent}{\footnotesize\textbf{[ado L#1]}}}

\newcommand{\R}{\mathbb{R}}
\newcommand{\E}{\mathbb{E}}
\newcommand{\1}{\mathbb{1}}
\newcommand{\Smpl}{\mathcal{S}}
\newcommand{\thetad}{\widehat{\bm{\theta}}_d}
\newcommand{\Ngt}{N_{g,t}}
\newcommand{\Pdca}{P^{\!W}_{\mathcal{F}_d}}
\newcommand{\Mdca}{M^{\!W}_{\mathcal{F}_d}}

\title{\textbf{Cluster-by-time HDFE in \texttt{did\_multiplegt\_dyn}} \\[2pt]
       \large Adding a focused option for FEs that are constant across groups within a time period}
\author{Worked example on \texttt{data\_test.csv}: \texttt{ac\_uq\_id\#monthyear}\\
        \small grid-cell panel where each grid (\texttt{unique\_small\_grid\_id}) belongs to a single \texttt{ac\_uq\_id}}
\date{}

\begin{document}
\maketitle

\section{The structure we want to absorb}

\begin{tcolorbox}[colback=codebg, colframe=navy, boxrule=0.6pt, arc=1pt]
\textbf{Dataset (\code{data\_test.csv}).}
Panel at the (group, time) level:
\begin{itemize}[leftmargin=*, itemsep=0pt]
\item \code{unique\_small\_grid\_id} \,---\, group id $g$,
\item \code{monthyear} \,---\, time id $t$,
\item \code{ac\_uq\_id} \,---\, a coarser group attribute that maps each
      \code{unique\_small\_grid\_id} to a single value for all $t$
      (\emph{time-invariant at the group level}).
\end{itemize}
\textbf{Goal.} Absorb \code{ac\_uq\_id\#monthyear} fixed effects \,---\, i.e.\
let each (cluster, time) cell have its own intercept \,---\, alongside the
time FE that \texttt{did\_multiplegt\_dyn} already partials out.
\end{tcolorbox}

\noindent
The defining feature of this FE: \emph{the absorbing variable is constant
over groups inside the cluster and over the time periods we want to allow it
to vary.} Formally,
\begin{equation}
a_g \;\equiv\; \mathrm{ac\_uq\_id}_g \in \{1,\dots,A\}\quad\text{is time-invariant,}\quad\text{and}\quad
a_g\times t \in \{1,\dots,A\}\times\{1,\dots,T\}
\label{eq:struct}
\end{equation}
takes one value per group at every $t$. This is the same structural
property that \texttt{trends\_nonparam} already exploits with one variable;
the proposed option generalizes it to multiple such variables and routes the
demeaning through \texttt{reghdfe}/\texttt{hdfe}.

\section{Why this case deserves its own option}

\begin{enumerate}[leftmargin=*, itemsep=2pt]
\item \textbf{The variance derivation in dCdH (2024)} is built for the
projection of $\Delta Y$ and $\Delta X_k$ onto time FEs (optionally
$\times$ a single time-invariant stratum). Cluster$\times$time FEs of the
form $a_g\times t$ are a \emph{linear extension} of that subspace and the
companion-paper algebra goes through verbatim.
\item \textbf{The companion paper's Section 1.4} already discusses
\texttt{trends\_nonparam}. We are not opening a new derivation \,---\,
we are saying ``allow multiple such strata, and use \texttt{reghdfe} as the
implementation back-end.''
\item \textbf{Performance:} a one-way demean by $(a_g, t, d)$ is the same
\code{bys ... : gegen} pattern already in the code at L931--934. Two-way
\,---\, e.g.\ \code{absorb(state\#year industry\#year)} \,---\, is what
forces us to call \texttt{reghdfe}.
\item \textbf{Restricted scope keeps validation simple.} The option does
\emph{not} accept slope FEs, time-varying categoricals, or unit FEs at a
finer level than the group. Only ``time-invariant group attribute $\times$
time.''
\end{enumerate}

\section{Proposed syntax}

\begin{lstlisting}
did_multiplegt_dyn Y G T D, effects(L) [other options] ///
                  cluster_time_fe(varlist)
\end{lstlisting}

\noindent
Semantics:
\begin{itemize}[leftmargin=*, itemsep=2pt]
\item Each variable $a^{(1)}, \dots, a^{(P)}$ in the \code{varlist} must be
      time-invariant at the group level (checked at runtime).
\item For each $a^{(p)}$, the option adds the FE set
      $\{\1\{a^{(p)}_g = a_0\}\cdot \1\{t = t_0\}\}_{a_0, t_0}$ to the
      projector used to residualize $\Delta Y$ and each $\Delta X_k$,
      \emph{separately per baseline level} $d$.
\item Time FE and (optionally) \code{trends\_nonparam} stay where they are;
      the new option just appends more dummies to $\mathcal{F}_d$.
\item If the user passes a variable that is not time-invariant at the group
      level, the command exits with a descriptive error.
\end{itemize}

\medskip\noindent\textbf{Concrete call on \code{data\_test.csv}}:
\begin{lstlisting}
import delimited using "data_test.csv", clear
did_multiplegt_dyn Y unique_small_grid_id monthyear D, ///
    effects(4) placebo(2) controls(X1 X2)             ///
    cluster_time_fe(ac_uq_id)                         ///
    cluster(ac_uq_id)
\end{lstlisting}

% ============================================================================
\section[Stage 1 -- The one stage that actually changes]{Stage 1 \,---\, The one stage that actually changes}

Recall Stage 1 partials a projector $\Pdca$ out of $\Delta Y$ and each
$\Delta X_k$, then runs OLS on the residuals to recover $\thetad$.
\textbf{All we change is the projector.}

\subsection*{Augmented FE space}

Let $a^{(p)}_g$ be the $p$-th \code{cluster\_time\_fe} variable. The
projector's FE space becomes:
\begin{equation}
\mathcal{F}_d \;=\; \mathrm{span}\Bigl\{\,
\underbrace{\1\{t=t_0\}}_{\text{time FE}},\
\underbrace{\1\{s_g = s_0\}\cdot\1\{t=t_0\}}_{\text{(optional) trends\_nonparam}},\
\underbrace{\1\{a^{(p)}_g = a_0\}\cdot \1\{t=t_0\}}_{\text{cluster\_time\_fe, $p=1,\dots,P$}}\,\Bigr\},
\label{eq:Fd}
\end{equation}
restricted to the control sample $\Smpl_d$
(not-yet-switchers with baseline $D_{g,1}=d$ and all $\Delta X_k$
non-missing). The weighted projector and its annihilator,
\[
\Pdca \;=\; \bm{Z}_d (\bm{Z}_d^{\top}\bm{W}_d \bm{Z}_d)^{+}\,\bm{Z}_d^{\top}\bm{W}_d,
\qquad
\Mdca \;=\; \bm{I} - \Pdca,
\]
project onto $\mathcal{F}_d$ in the weighted inner product
$\bm{W}_d = \mathrm{diag}(\Ngt)$.

\subsection*{Where the old Stage 1 lines change}

\begin{tcolorbox}[colback=codebg, colframe=ok, boxrule=0.6pt, arc=1pt]
\textbf{The math is identical to the existing Stage 1.} Replace every
appearance of ``demean by time FE within $\Smpl_d$'' with ``demean by
$\mathcal{F}_d$ within $\Smpl_d$.''
\end{tcolorbox}

\medskip\noindent
Concretely, the four substitutions:

\medskip
\noindent\textbf{(1) Residualized FD of each control \Adoref{941}.}
\begin{equation}
\widetilde{\Delta X}^{(\mathrm{new})}_{k,g,t} \;=\; \sqrt{\Ngt}\,\bigl[\Mdca\,\Delta X_k\bigr]_{g,t}
\quad\text{instead of}\quad
\widetilde{\Delta X}^{(\mathrm{old})}_{k,g,t} \;=\; \sqrt{\Ngt}\,\bigl(\Delta X_k - \overline{\Delta X}_k(d,t,s)\bigr).
\label{eq:resXnew}
\end{equation}

\medskip
\noindent\textbf{(2) Outcome FD residual \Adoref{955}.}
\begin{equation}
\widetilde{\Delta Y}^{(\mathrm{new})}_{g,t} \;=\; \sqrt{\Ngt}\,\bigl[\Mdca\,\Delta Y\bigr]_{g,t}
\quad\text{instead of}\quad
\widetilde{\Delta Y}^{(\mathrm{old})}_{g,t} \;=\; \sqrt{\Ngt}\,\Delta Y_{g,t}.
\label{eq:resYnew}
\end{equation}
(The old code did not partial time FE out of $\Delta Y$ here because
\texttt{matrix accum} computed $X^{\top}Y$ on raw $\Delta Y$ and the FE
absorption happened implicitly via the residualized $X$. With multi-way
absorb, FWL requires both sides to be demeaned by the same projector.)

\medskip
\noindent\textbf{(3) Per-control score \Adoref{958}.}
\begin{equation}
\mathrm{prod\_X}_{k,g,t}^{(\mathrm{new})} \;=\; \Ngt\cdot \bigl[\Mdca\,\Delta X_k\bigr]_{g,t},
\label{eq:prodXnew}
\end{equation}
mechanically replacing the same product against the old residual.

\medskip
\noindent\textbf{(4) Fitted FE-only mean of $\Delta Y$ \Adoref{984--996}.}
\begin{equation}
\widehat{E}^{(d,\mathrm{new})}_{g,t} \;=\; \bigl[\Pdca\,\Delta Y\bigr]_{g,t}
\quad\text{(via \texttt{reghdfe predict} \,---\, see code patch below).}
\label{eq:Ehatnew}
\end{equation}

\medskip\noindent
\textbf{What does NOT change in Stage 1.} After we compute the four
quantities above with the new projector:
\begin{itemize}[leftmargin=*, itemsep=2pt]
  \item \code{matrix accum overall\_XX = diff\_y\_wXX `mycontrols\_XX'}
        \Adoref{1018} works as before.
  \item $\thetad = \bm{M}_d^{-1}\bm{q}_d$ \Adoref{1034}: same formula.
  \item $\mathrm{Den}_d^{-1} = \bm{M}_d^{-1}\cdot N_d^{\mathrm{tot}}\cdot G$
        \Adoref{1052}: same formula.
  \item Singular-cell guard \Adoref{1042--1090}: same logic; just check
        $\mathrm{rank}(\bm{Z}_d) = \texttt{e(df\_a)}$ from \texttt{reghdfe}.
\end{itemize}

% ============================================================================
\section[Stages 2 and 3 -- No math change, just plug in the new theta]{Stages 2 and 3 \,---\, No math change, just plug in the new $\thetad$}

\begin{tcolorbox}[colback=codebg, colframe=ok, boxrule=0.6pt, arc=1pt]
\textbf{Zero code changes downstream.} Stages 2 and 3 consume
\code{coefs\_sq\_`d'\_XX}, \code{inv\_Denom\_`d'\_XX},
\code{prod\_X*\_Ngt\_XX}, and \code{E\_y\_hat\_gt\_int\_`d'\_XX}.\;
Each is rebuilt from the augmented projector \eqref{eq:Fd}, so the
downstream blocks at L4582--4689 (effects) and L5145--5184 (placebos), and
the variance correction at L4905--4954, automatically use the new HDFE
structure. No line in those blocks needs to be touched.
\end{tcolorbox}

\noindent
For reference, the only points where the Stage-1 outputs are read:
\begin{itemize}[leftmargin=*, itemsep=2pt]
\item L4679 \,---\, $\Delta Y^{(\ell)} \mathrel{-}{=} \thetad[k]\cdot \Delta X^{(\ell)}_k$.
\item L5177 \,---\, same line for the placebo long-difference.
\item L4647, L4664 \,---\, $(\Delta Y - \widehat{E}^{(d)}_{g,t})$ inside the score.
\item L4929 \,---\, $[\mathrm{Den}_d^{-1}\cdot \mathrm{in\_sum}_d]_{j,k}$.
\item L4938 \,---\, subtract $\thetad[j]$ from the influence-function block.
\item L4912--4946 \,---\, $\mathrm{prod\_X}_{k,g,t}$ multiplied inside $\mathrm{in\_sum}$.
\end{itemize}
Every consumer is satisfied as long as the Stage-1 outputs reflect the
augmented projection.

% ============================================================================
\section{Stata patch using \texttt{reghdfe}/\texttt{hdfe}}

Three pieces of patch: a syntax addition, the demean step that replaces
L919--961, and the FE-only fitted value that replaces L984--996.

\subsection*{Piece A \,---\, Syntax and dependency check (L53, L66)}

\begin{lstlisting}
* L53 -- add the option:
syntax varlist(min=4 max=4 numeric) [if] [in] [, ///
    /* ... existing options ... */              ///
    cluster_time_fe(varlist)]

* After the gtools check at L65-73, add:
if "`cluster_time_fe'" != "" {
    qui cap which reghdfe
    if _rc {
        di as error "cluster_time_fe() requires reghdfe. Run: ssc install reghdfe"
        exit 198
    }
    qui cap which hdfe
    if _rc {
        di as error "cluster_time_fe() also requires hdfe. Run: ssc install hdfe"
        exit 198
    }

    * Validate: each variable must be time-invariant at the group level.
    foreach v of varlist `cluster_time_fe' {
        tempvar sd_v
        bys `2': egen `sd_v' = sd(`v')
        sum `sd_v'
        if r(max) > 0 {
            di as error "cluster_time_fe() variable `v' is not time-invariant at the group level."
            exit 198
        }
        drop `sd_v'
    }
}
\end{lstlisting}

\subsection*{Piece B \,---\, Replace L919--961 (demean step)}

\begin{lstlisting}
* Build the FE set that the projector M_F_d will absorb.
local fe_set "time_XX"
if "`trends_nonparam'" != "" {
    capture drop trends_nonparam_temp_XX
    gegen trends_nonparam_temp_XX = group(`trends_nonparam')
    local fe_set "`fe_set' i.trends_nonparam_temp_XX#i.time_XX"
}
* The new piece: each cluster_time_fe variable, interacted with time.
foreach v of varlist `cluster_time_fe' {
    local fe_set "`fe_set' i.`v'#i.time_XX"
}

local mycontrols_XX ""
local count_controls = 0
foreach var of varlist `controls' {
    local ++count_controls
    capture drop resid_X`count_controls'_time_FE_XX
    capture drop prod_X`count_controls'_Ngt_XX
    gen resid_X`count_controls'_time_FE_XX = .
    gen prod_X`count_controls'_Ngt_XX = .
    local mycontrols_XX "`mycontrols_XX' resid_X`count_controls'_time_FE_XX"
}
capture drop diff_y_wXX
gen diff_y_wXX = .

levelsof d_sq_int_XX, local(levels_d_sq_XX)

foreach l of local levels_d_sq_XX {

    * Sample for baseline d = `l': not-yet-switcher controls.
    tempvar smp
    gen byte `smp' = (ever_change_d_XX == 0       ///
                    & diff_y_XX != .              ///
                    & fd_X_all_non_missing_XX==1  ///
                    & d_sq_int_XX == `l'          ///
                    & time_XX < F_g_XX)

    * One-shot multi-FE demean of DY and every DX_k.
    local dvars "diff_y_XX"
    forvalues k = 1/`count_controls' {
        local dvars "`dvars' diff_X`k'_XX"
    }
    qui hdfe `dvars' [aw=N_gt_XX] if `smp', ///
        absorb(`fe_set') generate(_resXX_)

    * Pack the residuals where the rest of the code expects them.
    replace diff_y_wXX = sqrt(N_gt_XX) * _resXX_diff_y_XX ///
        if d_sq_int_XX == `l' & `smp' == 1
    forvalues k = 1/`count_controls' {
        replace resid_X`k'_time_FE_XX = sqrt(N_gt_XX) * _resXX_diff_X`k'_XX ///
            if d_sq_int_XX == `l' & `smp' == 1
        replace resid_X`k'_time_FE_XX = 0 if missing(resid_X`k'_time_FE_XX) ///
            & d_sq_int_XX == `l'
        replace prod_X`k'_Ngt_XX = sqrt(N_gt_XX) * resid_X`k'_time_FE_XX ///
            if d_sq_int_XX == `l'
        replace prod_X`k'_Ngt_XX = 0 if missing(prod_X`k'_Ngt_XX) ///
            & d_sq_int_XX == `l'
    }
    drop _resXX_* `smp'
}
\end{lstlisting}

\subsection*{Piece C \,---\, Replace L984--996 (FE-only fitted value)}

\begin{lstlisting}
foreach l of local levels_d_sq_XX {

    capture drop E_y_hat_gt_int_`l'_XX

    * Build the same fitted value as before, but with the full FE set.
    * Note: we INCLUDE the controls in the regression so E_y_hat is the
    * full predicted value (FE part + X*theta_d), matching the existing
    * code at L996.
    local controlsXX ""
    forvalues k = 1/`count_controls' {
        local controlsXX "`controlsXX' diff_X`k'_XX"
    }

    cap reghdfe diff_y_XX `controlsXX' [aw=N_gt_XX]            ///
        if d_sq_int_XX == `l' & time_XX < F_g_XX,              ///
        absorb(`fe_set') noconstant residuals(_res_DY_l_XX)
    scalar rc_XX = _rc

    if rc_XX == 0 {
        * E_y_hat = DY - residual (matches what `predict` would do)
        gen E_y_hat_gt_int_`l'_XX = diff_y_XX - _res_DY_l_XX ///
            if d_sq_int_XX == `l' & time_XX < F_g_XX
        drop _res_DY_l_XX

        * Track absorbed-rank for the singularity guard at L1042.
        scalar df_a_`l'_XX = e(df_a)
        local useful_res_`l'_XX = e(N) - e(df_a) - `count_controls'
    }
    else {
        drop if d_sq_int_XX == `l'
        noi di "Baseline Treatment Level `l' dropped because of insufficient observations."
    }
}
\end{lstlisting}

\subsection*{The unchanged tail (L1018--1098)}

\begin{tcolorbox}[colback=codebg, colframe=ok, boxrule=0.6pt, arc=1pt]
The rest of the existing Stage-1 block \,---\, \code{matrix accum} at
L1018, $\thetad = \bm{M}_d^{-1}\bm{q}_d$ at L1034, the
$\mathrm{Den}_d^{-1}$ build at L1052, the singular-cell guard at L1042--1090,
and the dropping of cells at L1094--1096 \,---\, runs \emph{verbatim} on the
new residuals. No code change.
\end{tcolorbox}

% ============================================================================
\section{End-to-end example on \texttt{data\_test.csv}}

\begin{lstlisting}
* Path on the user's machine:
*   /Users/anzony.quisperojas/Documents/GitHub/did_multiplegt_dyn/Stata/data_test.csv

import delimited using "data_test.csv", clear

* Confirm the cluster-time FE candidate is time-invariant per group:
bys unique_small_grid_id: egen sd_ac = sd(ac_uq_id)
assert sd_ac == 0 | missing(sd_ac)
drop sd_ac

* Run the patched command:
did_multiplegt_dyn Y unique_small_grid_id monthyear D, ///
    effects(4) placebo(2)                             ///
    controls(X1 X2)                                   ///
    cluster_time_fe(ac_uq_id)                         ///
    cluster(ac_uq_id)
\end{lstlisting}

\noindent\textbf{What this estimates.} For each baseline-treatment level $d$:
\begin{equation}
\thetad \;=\; \arg\min_{\bm{\theta},\,\alpha^{(d)}}\;
\sum_{(g,t)\in\Smpl_d}\Ngt\Bigl(\Delta Y_{g,t} - \sum_{k}\theta_k\,\Delta X_{k,g,t} - \alpha^{(d)}_{a_g,\,t}\Bigr)^{2},
\label{eq:fwl}
\end{equation}
where $\alpha^{(d)}_{a,t}$ is an \code{ac\_uq\_id}-by-\code{monthyear} fixed
effect, allowed to differ by baseline cell. The event-study estimators
then use $Y_{g,t} - Y_{g,t-\ell} - (\Delta X^{(\ell)}_{g,t})^{\!\top}\thetad$
exactly as today, and the variance correction $U^{\mathrm{var},X}$ uses the
$\mathrm{Den}_d^{-1}$ built off this projection.

% ============================================================================
\section{Sanity check: equivalence with \texttt{trends\_nonparam}}

A regression test before merging: passing a single $a^{(1)}$ via
\code{cluster\_time\_fe(ac\_uq\_id)} should give numerically identical event
study coefficients to the current command run with
\code{trends\_nonparam(ac\_uq\_id)}, because the FE space \eqref{eq:Fd}
collapses to what \texttt{trends\_nonparam} already builds at L991. Run on
\code{data\_test.csv}:

\begin{lstlisting}
* OLD path:
did_multiplegt_dyn Y unique_small_grid_id monthyear D, ///
    effects(4) controls(X1 X2) trends_nonparam(ac_uq_id) ///
    cluster(ac_uq_id)
matrix b_old = e(b)
matrix V_old = e(V)

* NEW path:
did_multiplegt_dyn Y unique_small_grid_id monthyear D, ///
    effects(4) controls(X1 X2) cluster_time_fe(ac_uq_id) ///
    cluster(ac_uq_id)
matrix b_new = e(b)
matrix V_new = e(V)

* Compare element-wise:
matrix diff_b = b_old - b_new
matrix diff_V = V_old - V_new
matrix list diff_b
matrix list diff_V
\end{lstlisting}

\noindent
If the equivalence test passes (max abs diff below, say, $10^{-8}$), we have
a strong correctness signal for the single-FE case. Extending to multi-FE
(\code{cluster\_time\_fe(ac\_uq\_id another\_cluster)}) then has no
analytic counterpart, but the FWL argument and the per-baseline-cell
structure carry through.

% ============================================================================
\section{Edge cases}

\begin{enumerate}[leftmargin=*, itemsep=2pt]
\item \textbf{Group nested inside cluster.} If
      \code{unique\_small\_grid\_id} is finer than \code{ac\_uq\_id},
      \code{ac\_uq\_id} is the cluster. If they are 1-to-1 (each group is
      its own cluster), \code{ac\_uq\_id\#monthyear} reduces to the (g,t)
      cell which has only one obs \,---\, the regression is rank-deficient,
      and the singularity guard fires. The error message should hint at this.
\item \textbf{Singleton (cluster, year) cells.} \texttt{reghdfe} drops
      singletons by default. \code{nosingletons} is not what we want here
      \,---\, we should keep singletons but expect them to contribute zero
      after demeaning. Use \texttt{reghdfe}'s \code{keepsingletons} option;
      track \code{e(num\_singletons)} for diagnostics.
\item \textbf{Missing values in \code{cluster\_time\_fe} variables.}
      Treat as a separate level (\code{nomissing} off in \texttt{reghdfe}),
      or drop the observation. I recommend drop, with a one-line warning.
\item \textbf{Interaction with \texttt{trends\_nonparam}.} The two options
      compose: the FE set just stacks. \texttt{reghdfe} handles the
      collinearity between $a_g\times t$ terms and a $s_g\times t$ stratum
      that overlaps with $a_g$ via its rank-deficiency detector.
\item \textbf{Bootstrap and \texttt{by\_path()}.} Propagate
      \code{cluster\_time\_fe(`cluster\_time\_fe')} into the inner-call
      locals at L1337, L2030, L2033 \,---\, same as the existing options.
\end{enumerate}

% ============================================================================
\section{Summary}

\begin{tcolorbox}[colback=codebg, colframe=navy, boxrule=0.6pt, arc=1pt]
\textbf{The proposal in one paragraph.}\;
Add an option \code{cluster\_time\_fe(varlist)} that takes time-invariant
group attributes. Each is interacted with time and appended to the FE space
$\mathcal{F}_d$ used to residualize $\Delta Y$ and $\Delta X_k$ inside the
control sample for each baseline level $d$. The residualization is performed
by \code{hdfe} (one call per baseline cell, demeaning $\Delta Y$ and every
$\Delta X_k$ at once). The FE-only fitted value
$\widehat{E}^{(d)}_{g,t}$ is rebuilt by \code{reghdfe ..., absorb(\dots)
residuals()} and stored under the existing name. \emph{Nothing else changes}:
\code{matrix accum}, $\thetad$, $\mathrm{Den}_d^{-1}$, the long-difference
residualization at L4679, and the variance correction at L4905--4954 all
consume the rebuilt Stage-1 objects and produce correct event-study
estimates and variances by Frisch--Waugh--Lovell and the existing
companion-paper derivation.
\end{tcolorbox}

\end{document}
